%-*-latex-*-
%-*-latex-*-
\newcommand\COURSE{ciss350}
\newcommand\ASSESSMENT{q10603}
\newcommand\ASSESSMENTTYPE{Quiz}
\newcommand\POINTS{\textwhite{xxx/xxx}}

\input{myquizpreamble}
\input{yliow}
\input{\COURSE}
%-*-latex-*-

%-*-latex-*-

%-*-latex-*-
\renewcommand\TITLE{\ASSESSMENTTYPE \ \ASSESSMENT}

\renewcommand\EMAIL{}
\renewcommand\AUTHOR{eagraham1@cougars.ccis.edu}

\textwidth=6in
\begin{document}
\topmatterthree

For each of the following questions,
evaluate the expression written in either prefix/Polish
notation or postfix/Reverse Polish notation, or write ERROR if the expression is invalid.
You only need to write down the answer. You need not explain.

%------------------------------------------------------------------------------
\nextq
Polish notation:
\verb!* + 3 5 - 7 2!
\\
\ANSWER
\begin{answercode}
40
\end{answercode}

%------------------------------------------------------------------------------
\nextq
Polish notation: \verb!+ / 8 2 - 4 * 7 2!
\\
\ANSWER
\begin{answercode}
-6
\end{answercode}

%------------------------------------------------------------------------------
\nextq
Polish notation: \verb!* 8 * 2 + 3 - 2 / 7 2!
\\
\ANSWER
\begin{answercode}
32
\end{answercode}

%------------------------------------------------------------------------------
\nextq
Polish notation: \verb!/ + - * 2 3 7 4 2!
\\
\ANSWER
\begin{answercode}
0
\end{answercode}

%------------------------------------------------------------------------------
\nextq
Reverse Polish notation: \verb!2 3 4 * 5 * +!
\\
\ANSWER
\begin{answercode}
62
\end{answercode}

%------------------------------------------------------------------------------
\nextq
Reverse Polish notation: \verb!5 7 2 + 6 -!
\\
\ANSWER
\begin{answercode}
ERROR
\end{answercode}

%------------------------------------------------------------------------------
\nextq
Complete the following function that returns \verb!true! if the integer
parameter \verb!n! is a palindrome.
An integer is a palindrome if it's the same if the read it
left-to-right or right-to-left.
For instance \verb!0!, \verb!12321!, and \verb!123321! are palindromes.
Reading \verb!12321! left-to-right, you see \verb!1, 2, 3, 2, 1!,
which is the same if you read it right-to-left.
But \verb!-1! and \verb!1232! are not palindromes.
Note that reading \verb!-1! left-to-right gives you \verb!-, 1!
and reading it right-to-left gives you \verb!1, -!.
\\
\ANSWER
\begin{answercode}
bool is_palindrome(int n)
{
  if (n < 0) return false;

    int original = n;
    int reversed = 0;

    while (n > 0)
    {
        int digit = n % 10;
        reversed = reversed * 10 + digit;
        n /= 10;
    }

    return original == reversed;
}
\end{answercode}
(HINT: You can use a stack.
But you don't have to.
There are basically two ways to solve this problem,
you either look at the digits of $n$, or you form the reverse of $n$ with the
digits. There are variations of the above ideas.)

%------------------------------------------------------------------------------
%-*-latex-*-
\newpage
\input{instructions}
\end{document}

